\section{Localization and SLAM as Graphical Models}\label{sec:background}

\subsection{Dynamic Bayesian network}

A robot moves in the plane with pose $\x_t = (x_t, y_t, \theta_t)$ at time $t$. Between two time steps it receives
a control $\uu_t$ (a commanded motion) and then an observation $\z_t$ (ranges and bearings to landmarks). The map
$\m = (\m_1, \dots, \m_K)$ holds the positions of $K$ landmarks and does not change. These variables form a
dynamic Bayesian network \cite[ch.~6 and 15]{koller2009probabilistic}: each pose depends on the previous pose and
the control, and each observation depends on the current pose and the map (\cref{fig:dbn}). The joint distribution
factorizes accordingly:
\begin{equation}
  p(\x_{0:T}, \m, \z_{1:T} \mid \uu_{1:T}) = p(\x_0)\, p(\m) \prod_{t=1}^{T} p(\x_t \mid \x_{t-1}, \uu_t)\,
  p(\z_t \mid \x_t, \m).
  \label{eq:joint}
\end{equation}
The motion model $p(\x_t \mid \x_{t-1}, \uu_t)$ and the measurement model $p(\z_t \mid \x_t, \m)$ are the two
conditional probability distributions of the network, repeated at every time slice.

\begin{figure}[htbp]
  \centering
  \setlength{\unitlength}{1mm}
  \begin{picture}(110,42)
    % poses
    \put(10,20){\circle{10}}\put(10,20){\makebox(0,0){$\x_{t-2}$}}
    \put(40,20){\circle{10}}\put(40,20){\makebox(0,0){$\x_{t-1}$}}
    \put(70,20){\circle{10}}\put(70,20){\makebox(0,0){$\x_t$}}
    \put(15,20){\vector(1,0){20}}\put(45,20){\vector(1,0){20}}
    % controls
    \put(40,37){\circle{8}}\put(40,37){\makebox(0,0){$\uu_{t-1}$}}
    \put(70,37){\circle{8}}\put(70,37){\makebox(0,0){$\uu_t$}}
    \put(40,33){\vector(0,-1){8}}\put(70,33){\vector(0,-1){8}}
    % observations
    \put(40,3){\circle{8}}\put(40,3){\makebox(0,0){$\z_{t-1}$}}
    \put(70,3){\circle{8}}\put(70,3){\makebox(0,0){$\z_t$}}
    \put(40,15){\vector(0,-1){8}}\put(70,15){\vector(0,-1){8}}
    % map
    \put(100,20){\circle{10}}\put(100,20){\makebox(0,0){$\m$}}
    \put(95.5,18){\vector(-2,-1){21.5}}\put(95.5,17.5){\vector(-4,-1){52}}
  \end{picture}
  \caption{Two time slices of the dynamic Bayesian network for SLAM. Controls $\uu$ and observations $\z$ are
    observed; poses $\x$ and the map $\m$ are hidden. In localization the map is known.}
  \label{fig:dbn}
\end{figure}

In \emph{localization} the map is known and the task is to track the pose. Filtering means computing the belief
$p(\x_t \mid \z_{1:t}, \uu_{1:t})$ after each observation. Because $\x_t$ separates the past from the future in the
network, the belief satisfies the recursion
\begin{equation}
  p(\x_t \mid \z_{1:t}, \uu_{1:t}) \;\propto\; p(\z_t \mid \x_t)
  \int p(\x_t \mid \x_{t-1}, \uu_t)\, p(\x_{t-1} \mid \z_{1:t-1}, \uu_{1:t-1})\, d\x_{t-1},
  \label{eq:bayes-filter}
\end{equation}
a prediction through the motion model followed by a correction with the likelihood of the new observation
\cite{thrun2005probabilistic}. This is the forward pass of an unrolled network: sum-product message passing along
the chain of poses.

In \emph{SLAM} the map is unknown as well, and the target is the joint posterior $p(\x_{1:t}, \m \mid \z_{1:t},
\uu_{1:t})$. Throughout this report the correspondence between observations and landmarks is known: each
observation carries the identity of the landmark it measures.

\subsection{Rao-Blackwellization}

The map has a structure that the network makes visible. Write the map as one node per landmark. An observation of
landmark $k$ at time $s$ has exactly two parents, $\x_s$ and $\m_k$, because the correspondence is known. A path from
$\m_j$ to $\m_k$ therefore has to leave $\m_j$ through one of its observations and reach that observation's pose
$\x_s$, where both edges of the path point away from $\x_s$. If the whole trajectory $\x_{1:t}$ is given, $\x_s$
blocks the path. The observations on the path are observed colliders, but every such path continues through a given
pose, so all paths are blocked: the landmarks are d-separated given the trajectory and the observations, and
therefore conditionally independent \cite{koller2009probabilistic}. The posterior factorizes as
\begin{equation}
  p(\x_{1:t}, \m \mid \z_{1:t}, \uu_{1:t}) = p(\x_{1:t} \mid \z_{1:t}, \uu_{1:t})
  \prod_{k=1}^{K} p(\m_k \mid \x_{1:t}, \z_{1:t}).
  \label{eq:rao-blackwell}
\end{equation}
The independence holds only given the full trajectory; given only the current pose, the landmarks are correlated
through the uncertain path that observed them.

Rao-Blackwellization \cite{doucet2000rao} uses this factorization to split the problem. The trajectory, the hard
part with nonlinear and possibly multimodal uncertainty, is represented by samples. Given a sampled trajectory each
factor $p(\m_k \mid \x_{1:t}, \z_{1:t})$ is a small estimation problem over two variables, which FastSLAM solves with
one extended Kalman filter per landmark \cite{montemerlo2002fastslam}. Computing part of the posterior analytically
instead of sampling it can only reduce the variance of the estimates: by the law of total variance,
$\operatorname{Var}[\mathbb{E}[f \mid \x_{1:t}]] \le \operatorname{Var}[f]$. It also keeps the cost low. Each
particle holds $K$ Gaussians of dimension two, instead of one Gaussian of dimension $2K + 3$, and an observation
updates only the landmark it measures.

\subsection{From discrete to continuous inference}

For discrete variables the recursion \eqref{eq:bayes-filter} is the forward algorithm of a hidden Markov model: the
belief is a table and the integral is a sum, computed exactly. For linear models with Gaussian noise the belief
stays Gaussian and the recursion is the Kalman filter. Neither applies here. The measurement model, a range and a
bearing to a landmark, is nonlinear, and after a uniform start the belief is spread over the whole world and has
many modes. Linearizing, as the extended Kalman filter does, keeps a single mode and fails when the belief is far
from Gaussian.

What carries over from the discrete case is the structure: the factorization \eqref{eq:joint}, the conditional
independences used by \eqref{eq:bayes-filter} and \eqref{eq:rao-blackwell}, and the message-passing view of
filtering. What changes is how a message is represented. Sampling methods replace exact summation with sums over
samples: likelihood weighting in a Bayesian network weights each sample by the likelihood of the evidence, which is
importance sampling with the prior as the proposal \cite[ch.~12]{koller2009probabilistic}. The particle filter
applies this idea at every time slice of the dynamic network and adds resampling so that the samples stay useful
over time \cite[ch.~15]{koller2009probabilistic}.

\subsection{Learning the parameters}\label{sec:learning}

The two conditional distributions above are given, but each carries a noise level that somebody has to choose.
Collect those numbers in $\theta$: here four standard deviations, for the forward motion, the turn, the measured
range and the measured bearing. Estimating them from data is \emph{parameter} learning rather than structure
learning, since the network of \cref{fig:dbn} is fixed and only the numbers inside two of its conditional
distributions are unknown. It is parameter learning \emph{with incomplete data}, because the trajectory the
observations depend on is never observed \cite[ch.~19]{koller2009probabilistic}: the data are the controls and the
observations of a recorded run, and nothing else.

A single $\theta$ covers the whole run because a dynamic network is a template model. The same conditional
distribution is instantiated at every time slice, so its parameters are tied across slices
\cite[ch.~6]{koller2009probabilistic}. Every one of the $T$ steps is then evidence about the same four numbers,
and a longer trajectory is more data about $\theta$ rather than more unknowns.

The maximum-likelihood estimate is the $\theta$ under which the recorded observations are most probable,
\begin{equation}
  \hat{\theta} = \arg\max_{\theta}\; \log p(\z_{1:T} \mid \uu_{1:T}, \theta), \qquad
  p(\z_{1:T} \mid \uu_{1:T}, \theta) = \prod_{t=1}^{T} p(\z_t \mid \z_{1:t-1}, \uu_{1:t}, \theta).
  \label{eq:mle}
\end{equation}
Each factor integrates over all trajectories and has no closed form, but the particle filter estimates it as a
by-product. With $\w_{t-1}^{(i)}$ the normalized weights after step $t-1$ and particles propagated through the
motion model, the unnormalized weight of particle $i$ is
$\tilde{\w}_t^{(i)} = \w_{t-1}^{(i)}\, p(\z_t \mid \x_t^{(i)}, \theta)$ and
\begin{equation}
  p(\z_t \mid \z_{1:t-1}, \uu_{1:t}, \theta) \;\approx\; \sum_{i=1}^{N} \tilde{\w}_t^{(i)},
  \label{eq:evidence}
\end{equation}
which is the sum the filter forms anyway in order to normalize its weights \cite{doucet2011tutorial}. Adding the
logarithm of \eqref{eq:evidence} over the run estimates $\log p(\z_{1:T} \mid \theta)$. The estimate of the
likelihood is unbiased, so the estimate of its logarithm is biased downwards, and that bias grows as the particle
approximation degrades. The same machinery is therefore both the inference and the objective that learning
maximizes, which ties the quality of one to the other.

Two properties of \eqref{eq:mle} decide whether anything can be learned from it. The first is that its maximum
is in the right place. As the run grows the average log likelihood tends to an expectation that splits into a term
free of $\theta$ and a Kullback-Leibler divergence,
\begin{equation}
  \frac{1}{T} \log p(\z_{1:T} \mid \theta) \;\longrightarrow\;
  \mathbb{E}_{p^*}\!\left[\log p^*(\z)\right] - \mathrm{KL}\!\left(p^* \,\|\, p_{\theta}\right),
  \label{eq:kl}
\end{equation}
and the divergence is non-negative and zero only where the two distributions agree
\cite[ch.~17]{koller2009probabilistic}. The expected score is therefore largest at the parameters that generated
the data: maximizing the likelihood is minimizing the distance to the true model.

The second is how sharp that maximum is. Expanding the expected score around $\theta^*$ gives
\begin{equation}
  \mathbb{E}\!\left[\log p(\z_{1:T} \mid \theta)\right] \;\approx\; \text{const}
  - \tfrac{1}{2}\, T\, \mathcal{I}(\theta^*)\, (\theta - \theta^*)^2,
  \label{eq:fisher}
\end{equation}
with $\mathcal{I}$ the information a single step carries about $\theta$, and the variance of the estimate is
$1/(T \mathcal{I})$. A parameter the observations say nothing about has $\mathcal{I}$ near zero and a flat
likelihood: it is part of the model and still not identifiable from the run. Because the parameters are tied
across slices, the curvature grows with the number of steps, so a parameter that a short trajectory cannot
resolve can be resolved by a longer one.

The shape on either side of the maximum follows from the same expression. For one Gaussian residual with true
standard deviation $\sigma^*$, scored under an assumed $\sigma$,
\begin{equation}
  \mathbb{E}\!\left[\log p\right] = -\log \sigma - \frac{\sigma^{*2}}{2 \sigma^{2}} - \tfrac{1}{2} \log 2\pi,
  \label{eq:gaussian-score}
\end{equation}
whose derivative $-1/\sigma + \sigma^{*2}/\sigma^{3}$ vanishes at $\sigma = \sigma^*$. The two terms are not
symmetric: assuming too little noise is penalized like $\sigma^{-2}$, assuming too much only like $\log \sigma$.
Understating a noise level is therefore punished far more heavily than overstating it, which is a reason to set
such a value generously when it is set by hand.

Maximizing \eqref{eq:mle} by sweeping one component of $\theta$ at a time, with the others held fixed, is a profile
likelihood. A general method would instead iterate, alternating a step that reweights trajectories under the current
$\theta$ with a step that updates $\theta$ in closed form, which is expectation maximization
\cite[ch.~19]{koller2009probabilistic} applied to a state-space model. The sweep is used here because it also shows
the shape of the likelihood, and that shape is what decides whether a parameter can be identified at all.
